\section{The certificate method}\label{sec:certificates}

Every exclusion in this paper has the same shape: find, in the cone relation of a run, a linear-size acyclic set, possibly after
revealing a cheap set of steps, and conclude that the run cannot belong to a linear-time transposer. This section states the
method and the tools used to find acyclic sets.

\subsection{Genies and index codes}

\begin{lemma}[Genies and index codes]\label{lem:genie}
\cited{\cite[Corollary~O]{BoundedStack2026}, Lean-verified there (\lean{BoundedStackDepth.genie_cost},
\lean{BoundedStackDepth.genie_of_trace}, \lean{BoundedStackDepth.exists_indexCode}).}
\begin{enumerate}[label=(\alph*)]
\item For every $k\ge1$, $\beta\ge1$, $C\ge1$, $\eps>0$ and every trace with $k$ work tapes and $T\le C(L+1)$, there is a set $\Gamma$
of steps with $\beta\abs\Gamma\le\eps L$ such that every dependency path whose vertices after the first avoid $\Gamma$ has at most
$\mathrm{genieDepth}(k,\beta,C,\eps)$ vertices. The depth bound depends only on $k$, $\beta$, $C$ and~$\eps$.
\item If $M$ is oblivious at length $L$ with common trace $\tau$ and transposes, then for every set $\Gamma$ of steps the values
written at the steps of $\Gamma$ form an index code of $\wval\abs\Gamma$ bits for the cone relation avoiding $\Gamma$: two inputs
with the same values on $\Gamma$ that agree on the cone of $\Out(q)$ agree on the output bit at $\Out(q)$.
\end{enumerate}
\end{lemma}

The genie comes from the fourteenth paper's solution of the trajectory cut problem~\cite[Theorem~T]{BoundedStack2026}; its depth
$\mathrm{genieDepth}(k,\beta,C,\eps)$ is explicit, with $\log\log\mathrm{genieDepth}=O(k^2\log(k\beta C/\eps))$~\cite[Remark~9.4 and
\S10]{BoundedStack2026}.

\begin{lemma}[Acyclic sets bound codes]\label{lem:acyclic-code}
\formal{DichromaticCrux.lean}{storageCode_card_le_of_isAcyclicIn, indexCode_bits_ge_of_isAcyclicIn}
Let $N$ be a relation on a finite set $V$ and $A\subseteq V$ acyclic in~$N$. Every storage code for $N$ over an alphabet $\Sigma$ has
at most $\abs\Sigma^{\abs V-\abs A}$ words, and every index code for $N$ and the demand $\tpos$ has at least $\abs A$ bits, provided
$A$ is acyclic for $q\mapsto N(\tpos(q))$.
\end{lemma}

\begin{outline}
On a storage code a word is determined by its letters outside $A$: recover the letters in $A$ in an order in which each letter's
neighbourhood lies outside the not-yet-recovered part (\lean{acyclicIn_outside_injOn}). The fibres of an index code are storage
codes for the demand-shifted relation, and pigeonhole over the messages gives the bound.
\end{outline}

\subsection{Cruxes}

\begin{definition}[Cruxes]\label{def:crux}
\formal{MachineCrux.lean, EnvelopeCrux.lean}{SubpopCutCruxAt, MachineCutCruxAt, MachineCruxAt, subpopCutCruxAt_one_iff, machineCutCruxAt_zero_iff}
Let $M$ be a machine, $C\ge1$, $D\ge0$, $0<\alpha_0\le1$ and $\rho,\delta\ge0$. $M$ has \emph{subpopulation cut-crux rate $\rho$ at
density $\alpha_0$, cost $\delta$ and depth $D$} if for every $\eta>0$ and all large $d$, whenever $M$ is oblivious at length $L$
with halting time $T\le C(L+1)$ and exactly $L$ appends, there are a set $V_0$ of at least $\alpha_0L$ input positions and a set $B$
of steps with $\wval\abs B\le\delta\abs{V_0}$ such that every index code for the cone relation avoiding $\Gamma\cup B$, for any
$\Gamma$ that cuts all dependency paths to at most $D+1$ vertices, has at least $(\rho-\eta)\abs{V_0}$ bits (\lean{SubpopCutCruxAt}).
Two cases have their own names.
\begin{enumerate}[label=(\alph*)]
\item At density $1$ and cost $0$, $M$ has \emph{crux rate $\rho$ at depth $D$}. Then $V_0=[0,L)$ and $B=\emptyset$, since $\wval\ge2$:
every index code for the cone relation avoiding $\Gamma$ has at least $(\rho-\eta)L$ bits (\lean{MachineCruxAt},
\lean{machineCutCruxAt_zero_iff}).
\item At density $1$, $M$ has \emph{cut-crux rate $\rho$ at cost $\delta$ and depth $D$}. Then $V_0=[0,L)$: there is a set $B$ with
$\wval\abs B\le\delta L$ such that every index code for the cone relation avoiding $\Gamma\cup B$ has at least $(\rho-\eta)L$ bits
(\lean{MachineCutCruxAt}, \lean{subpopCutCruxAt_one_iff}).
\end{enumerate}
\end{definition}

\begin{theorem}[Cruxes exclude]\label{thm:crux-excludes}
\formal{EnvelopeCrux.lean, MachineCrux.lean}{subpop_cutcrux_excludes, machine_cutcrux_excludes, machine_crux_excludes}
Let $C\ge1$ and $\eps>0$, and put $D_\eps=\mathrm{genieDepth}(k,\wval,C,\eps)$. If $M$ has subpopulation cut-crux rate $\rho$ at
density $\alpha_0$, cost $\delta$ and depth $D_\eps$ with $\eps<\alpha_0(\rho-\delta)$, then for all large $d$ there is no $T$ with
$\mathrm{ObliviousLinTime}(M,d,C,T)$. In particular a cut-crux rate $\rho$ at cost $\delta$ with $\eps+\delta<\rho$, and a crux rate
$\rho>\eps$, at depth $D_\eps$ exclude.
\end{theorem}

\begin{outline}
Suppose $M$ transposes obliviously in time $C(L+1)$. Reveal the genie $\Gamma$ of \cref{lem:genie}(a), of cost $\eps L$ and depth
$D_\eps$, together with~$B$. By \cref{lem:genie}(b) the values on $\Gamma\cup B$ form an index code of at most $\eps L+\delta\abs{V_0}$
bits, so $(\rho-\delta-\eta)\abs{V_0}\le\eps L$, which contradicts $\abs{V_0}\ge\alpha_0L$ for small $\eta$. The two special cases
are $\alpha_0=1$, where $\eps<\alpha_0(\rho-\delta)$ reads $\eps+\delta<\rho$, and in addition $\delta=0$, where it reads $\eps<\rho$;
their direct forms are \lean{machine_cutcrux_excludes} and \lean{machine_crux_excludes}.
\end{outline}

Acyclic sets give cruxes. If, for every $\eta>0$ and all large $d$, every cone relation of depth $D$ in \cref{def:crux}(a) contains
an acyclic set of at least $(\rho-\eta)L$ input positions, then $M$ has crux rate $\rho$ at depth~$D$, by \cref{lem:acyclic-code}
(\lean{machineCruxAt_of_isAcyclicIn}, \lean{machineCrux_of_isAcyclicIn}, \lean{oblivious_excludes_of_isAcyclicIn}). The same holds for ball relations, which contain the
cone relations of depth~$D$ (\lean{machineCutCruxAt_of_ballAvoid_isAcyclicIn}).

Part~(b) of the next lemma is the form used by Theorems~X and~S: at a single size, revealing a set $B$ that leaves an acyclic set $U$ in
the cone relation costs at least $\abs U$ bits. Part~(a), used by the first forms, does the same for the ball relation of a bounded
radius, up to the genie.

\begin{lemma}[Acyclic sets cost cuts]\label{lem:acyclic-cost}
\formal{EnvelopeCrux.lean, TightConstants.lean}{obliviousLinTime_separator_card_ge_of_isAcyclicIn, obliviousLinTime_coneSystem_isAcyclicIn_card_le, obliviousLinTime_ballSeparator_card_ge}
Let $\mathrm{ObliviousLinTime}(M,d,C,T)$ with $d\ge1$, $C\ge1$ and $\eps>0$, put $D_\eps=\mathrm{genieDepth}(k,\wval,C,\eps)$, and let
$B$ be a set of steps.
\begin{enumerate}[label=(\alph*)]
\item For every set $U$ of input positions that is acyclic in the ball relation of radius $D_\eps$ avoiding $B$,
$\wval\abs B\ge\abs U-\eps L$.
\item For every set $U$ of input positions that is acyclic in the cone relation avoiding $B$, $\wval\abs B\ge\abs U$; this part needs
only $\mathrm{ObliviousLinTime}(M,d,C,T)$.
\end{enumerate}
\end{lemma}

\begin{outline}
(a) Reveal the genie of cost $\eps L$ together with $B$. The cone relation avoiding both is contained in the ball relation
(\lean{coneSystem_subset_ballAvoidSystem}), so $U$ is acyclic in it, and \cref{lem:genie}(b) with \cref{lem:acyclic-code} gives
$\wval(\abs\Gamma+\abs B)\ge\abs U$. (b) The values on $B$ alone already form an index code for the cone relation avoiding $B$,
which has no depth bound, and \cref{lem:acyclic-code} applies with no genie.
\end{outline}

\subsection{Heavy traces of arbitrary machines}

For a machine that is not oblivious, the index code of \cref{lem:genie}(b) is not available, but counting the fibre of a single
trace replaces it.

\begin{lemma}[Frozen-trace deficit]\label{lem:frozen}
\formal{SimulatorBridge.lean}{correctSubset_isAcyclicIn_coneSystem_deficit_le_logb, correctSubset_isAcyclicIn_coneSystem_deficit_le, frozenTrace_isAcyclicIn_coneSystem_deficit_le, frozenTrace_isAcyclicIn_ballAvoid_separator_card_ge}
Let $M$ be any machine, $d\ge1$, $\tau$ a trace of $M$ with halting time $T$, and $S$ a nonempty set of inputs that follow $\tau$ and
that $M$ transposes within time $c$. Then $T\le2cL$, $\tau$ appends exactly $L$ bits, and for every set $\Gamma$ of steps and every set
$U$ acyclic in the cone relation of $\tau$ avoiding $\Gamma$,
\[
  \abs U-\wval\abs\Gamma\ \le\ L-\log_2\abs S .
\]
In particular $\abs U-\wval\abs\Gamma\le\eta L$ if $\abs S\ge2^{(1-\eta)L}$; for instance, if $M$ transposes at size $d$ within time
$c$ and $\tau$ is $\eta$-heavy, $S$ can be the set of all inputs that follow $\tau$.
\end{lemma}

\begin{outline}
All inputs of $S$ share the dependency DAG of $\tau$, and any one of them gives the facts about $T$. Split $S$ according to the
$2^{\wval\abs\Gamma}$ possible values on $\Gamma$. Since the inputs of $S$ are transposed, the bit appended at $\Out(q)$ is input
bit $q$, so each part is a storage code for the cone relation and has at most $2^{L-\abs U}$ elements by \cref{lem:acyclic-code}.
Hence $\abs S\le2^{\wval\abs\Gamma+L-\abs U}$. Correctness is used only on~$S$.
\end{outline}

\subsection{Tools for finding acyclic sets}

Acyclic sets in cone relations are found by potentials. The following statements hold for any relation $N$ on a finite set; an arc
$q\to p$ means $p\in N(q)$.

\begin{lemma}[Envelopes and windows]\label{lem:window}
\formal{EnvelopeCrux.lean}{acyclicIn_of_lt_on_nbhd, upperEnvelope_isAcyclicIn, oneSidedEnvelope_isAcyclicIn, exists_isAcyclicIn_of_oneSidedEnvelope_maxDeg, window_isAcyclicIn_of_noFarArcs, exists_isAcyclicIn_of_window_maxDeg}
\begin{enumerate}[label=(\alph*)]
\item If a potential strictly decreases along every arc inside $S$, then $S$ is acyclic.
\item Let $\varphi,\psi$ be real potentials and $w$ a real number. The \emph{window} $\{q\in V_0:q\notin N(q),\ \varphi(q)+w<\psi(q)\}$
is acyclic if every arc $q\to p$ inside it is \emph{near}, $\psi(p)\le\varphi(q)+w$. If instead every vertex of the window is
incident to at most $\beta$ far arcs inside it, the window contains an acyclic set of at least a $1/(\beta+1)$ fraction of its
elements.
\end{enumerate}
\end{lemma}

\begin{outline}
(a) Inside $S$ an element of least potential has no out-arc into $S$. (b) A near arc $q\to p$ inside the window has
$\varphi(p)<\psi(p)-w\le\varphi(q)$, so $\varphi$ strictly decreases along it and (a) applies. With far arcs, remove the far neighbours
of the vertices one at a time; a greedy choice keeps a $1/(\beta+1)$ fraction. A two-sided bound $\abs{\psi(p)-\varphi(q)}\le w$ is
near at margin $2w$; with two-sided bounds the statements are \lean{window_isAcyclicIn_of_noFarArcs} and
\lean{exists_isAcyclicIn_of_window_maxDeg}. The upper-envelope form of (a), that $\{q\in V_0:\Psi(q)<\varphi(q)\}$ is acyclic when
$\varphi(r)\le\Psi(p)$ for every arc $r\to p$ inside $V_0$, is \lean{upperEnvelope_isAcyclicIn}.
\end{outline}

\begin{lemma}[Cut sets and block ranking]\label{lem:blockcut}
\formal{EnvelopeCrux.lean}{cutSet_isAcyclicIn, exists_isAcyclicIn_card_ge_of_blockCut, exists_isAcyclicIn_card_ge_of_uniformBlockCut, exists_isAcyclicIn_of_shiftBlockCutFamily_gen}
Assign to each $q\in V_0$ two disjoint sets $A(q),E(q)$ of blocks in $\{0,\dots,B-1\}$, and suppose every arc $q\to p$ inside $V_0$ has
$A(q)\cap E(p)\ne\emptyset$. For a ranking $\sigma$ of the blocks, the set of $q$ with $\sigma(A(q))<\sigma(E(q))$ (every block of $A(q)$
ranked below every block of $E(q)$) is acyclic, and some ranking makes it contain at least
$\sum_q\binom{\abs{A(q)}+\abs{E(q)}}{\abs{A(q)}}^{-1}$ elements, hence at least $\abs{V_0}/\binom{a+e}{a}$ elements when
$\abs{A(q)}\le a$ and $\abs{E(q)}\le e$ throughout.
\end{lemma}

\begin{outline}
Along an arc $q\to p$ inside the set, a block of $A(q)\cap E(p)$ shows that the largest rank in $A(p)$ is below the smallest rank in $E(p)$,
hence below that block's rank, hence at most the largest rank in $A(q)$; so the largest rank of $A$ strictly decreases along arcs, and
\cref{lem:window}(a) applies. For disjoint $A(q),E(q)$, exactly a fraction $\binom{\abs{A(q)}+\abs{E(q)}}{\abs{A(q)}}^{-1}$ of all rankings
put $A(q)$ below $E(q)$ (\lean{LeanMachine.BlockRanking.choose_mul_card_goodPerms_eq_factorial}), so the average size of the set over all rankings is the stated
sum. For a family of $m$ block assignments, of which $s$ elements in total exceed the sizes $a,e$ and at most $\kappa$ per assignment
have $A(q)\cap E(q)\ne\emptyset$, the assignment with the fewest exceptions gives an acyclic set of at least
$(\abs{V_0}-s/m-\kappa)/\binom{a+e}{a}$ elements (\lean{exists_isAcyclicIn_of_shiftBlockCutFamily_gen}).
\end{outline}

\begin{lemma}[Step-invariant windows]\label{lem:stepinv}
\formal{EnvelopeCrux.lean}{isAcyclicIn_ballAvoid_window_of_stepInvariant, isAcyclicIn_ballAvoid_window_of_constStepInvariant}
Let a trace, a set $B$ of steps, a depth $D$ and a set $V_0$ of input positions be given, with regions $R(0),\dots,R(D)$ of steps, a
\emph{sink} region, a step potential $\mu$, potentials $\varphi,\psi$ on input positions, reals $c_0,c_1$ and $\Delta\ge0$. Suppose:
the roots of $V_0$ lie in $R(D)$ with $\mu(\Out(q))\le\varphi(q)+c_0$; every edge from the sink that avoids $B$ stays in the sink,
and the sink reads no position of $V_0$; every edge $u\to v$ from $R(r+1)$ that avoids $B$ enters the sink or enters $R(r)$ with
$\mu(v)\le\mu(u)+\Delta$; and a read of $p\in V_0$ at a step $u\in R(r)$ has $\psi(p)\le\mu(u)+c_1$. Then
$\{q\in V_0:\varphi(q)+c_0+D\Delta+c_1<\psi(q)\}$ is acyclic in the ball relation of radius $D$ avoiding~$B$.
\end{lemma}

\begin{outline}
By induction along a dependency path from a root with at most $D$ edges that avoids $B$ after the root, its end lies in the sink or in
the region $R(D-\ell)$ after $\ell$ edges, with $\mu$ raised by at most $\ell\Delta\le D\Delta$
(\lean{depPath_oneSidedStepInvariant_with_sink}). Since the sink reads nothing in $V_0$, every arc $q\to p$ inside $V_0$ has
$\psi(p)\le\varphi(q)+c_0+D\Delta+c_1$: it is near, and \cref{lem:window}(b) applies. The strict inequality excludes loops, so no
sign condition on $c_0,c_1$ is needed. The condition $\Delta\ge0$ is needed: for $D=1$ and $\Delta=-1$, a root that reads its own
position is a loop inside the set.
\end{outline}

\Cref{lem:window,lem:blockcut,lem:stepinv} are statements about traces only. They are the tools behind the exclusions of the
schedule families in \cref{app:catalogue}. The first form of \cref{thm:pullback} (\lean{timePullback_blockCut_mais}) was a far-reaching use of \cref{lem:blockcut};
the present form replaces the uniform block cut by a single interval on which the block sets are constant.
